\begin{tikzpicture}[
    >=Latex,
    scale=1.3,
    every node/.style={font=\small},
]
    % --- Buildplate ---
    \fill[gray!15] (-2.0, -0.4) rectangle (3.5, 0);
    \draw[gray!60] (-2.0, -0.4) rectangle (3.5, 0);

    % --- Bauteil: schmaler Stamm + breiter Auswuchs nach rechts oben ---
    \fill[blue!10]
        (-1.0, 0) -- (-1.0, 1.8)
        .. controls (-0.8, 2.5) and (1.0, 2.6) ..
        (2.0, 2.4) --
        (2.5, 1.8) --
        (2.5, 1.2) --
        (1.0, 1.0) --
        (1.0, 0) -- cycle;
    \draw[thick, blue!60!black]
        (-1.0, 0) -- (-1.0, 1.8)
        .. controls (-0.8, 2.5) and (1.0, 2.6) ..
        (2.0, 2.4) --
        (2.5, 1.8) --
        (2.5, 1.2) --
        (1.0, 1.0) --
        (1.0, 0);
    \node[blue!60!black, font=\scriptsize]
        at (-1.5, 0.9) {Bauteil};

    % =========================================================
    % Eltern-Lauf: Mittelpunkt c und Sphären
    % =========================================================
    % Mittelpunkt c
    \fill[blue!70!black] (0, 0) circle (2.5pt);
    \node[blue!70!black, anchor=north, font=\scriptsize]
        at (0, -0.05) {$\vec{c}$ (Eltern-Lauf)};

    % S_{i-1} (letzte druckbare Schale)
    \draw[thick, blue!70!black]
        (0, 0) ++ (0:1.5) arc (0:180:1.5);
    \node[blue!70!black, font=\scriptsize, anchor=south west]
        at ($(0, 0) + (50:1.5)$) {$S_{i-1}$};

    % S_i (mit Überhang)
    \draw[thick, blue!70!black, dashed]
        (0, 0) ++ (0:1.9) arc (0:180:1.9);
    \node[blue!70!black, font=\scriptsize, anchor=south west]
        at ($(0, 0) + (35:1.9)$) {$S_i$};

    % Überhängender Bereich auf S_i (rot markiert, über dem Auswuchs)
    \draw[ultra thick, red!80!black]
        ($(0, 0) + (15:1.9)$) arc (15:55:1.9);
    \node[red!80!black, font=\tiny, anchor=west]
        at ($(0, 0) + (35:2.05)$) {Überhang};

    % =========================================================
    % Neues Sphärenzentrum c' und Kind-Lauf
    % =========================================================
    % c' liegt auf S_{i-1} am Schwerpunkt des Überhang-Bereichs
    \coordinate (cprime) at ($(0, 0) + (35:1.5)$);
    \fill[orange!85!black] (cprime) circle (2.5pt);
    \node[orange!85!black, anchor=west, font=\scriptsize, align=left]
        at ($(cprime) + (0.1, 0.0)$) {$\vec{c}'$\\\scriptsize (Kind-Lauf)};

    % Neue Sphäre S'_0 vom c' aus, kleiner Radius
    \draw[thick, orange!85!black]
        ($(cprime) + (-65:0.3)$) arc (-65:115:0.3);
    \draw[thick, orange!85!black]
        ($(cprime) + (-65:0.5)$) arc (-65:115:0.5);

    % Kennzeichnen: Wachstumsrichtung des Kind-Laufs (radial von c')
    \draw[->, very thick, orange!85!black]
        ($(cprime) + (35:0.65)$) -- ++(35:0.4);
    \node[orange!85!black, font=\scriptsize, anchor=west]
        at ($(cprime) + (35:1.1)$) {$\vec{G}'$};

    % =========================================================
    % Annotation des Spawn-Ereignisses
    % =========================================================
    \draw[->, thin, gray!60!black]
        ($(0, 0) + (35:1.95)$) .. controls (1.7, 2.4) and (2.0, 2.0) .. (cprime);
    \node[gray!60!black, font=\tiny, anchor=south west,
            align=left]
        at (1.5, 1.95) {neuer Mittelpunkt $\vec{c}'$\\auf $S_{i-1}$ am Schwerpunkt\\des Überhangs};

\end{tikzpicture}